\scanpage{44}

Variantes évidentes avec des variétés topologiques compactes.
On peut prendre pour $k$ un anneau \inserted{(\unsure{pas corps})}, lorsqu'on a une théorie de \unsure{cohomologie|coh. alg.} à coefficients dans $k$ \deleted{\unsure{qui satisfait}} \inserted{telle que} \deleted{\unsure{Poincaré}} $H^*(X)$ satisfait à Poincaré \deleted{\illegible{1}} (pas seulement mod torsion\ldots).

2) Considérons un autre foncteur de $\unsure{\mathcal{V}}$ en les anneaux bivariants \unsure{jaugés}, tel le foncteur $A^*(X)$ (anneau de Chow, \deleted{\unsure{gradué}} \unsure{jaugé} par le \uline{degré} des cycles de degré max), ou un foncteur quotient (tels les groupes de cycles mod équiv. numérique, homologique, $\unsure{\checkmark^{\tau}}$ algébrique, Picard, Albanese\ldots). On trouve un anneau bivariant \unsure{jaugé} simplicial, qui n'a plus aucune tendance à être spécial, et dont les composantes n'ont pas \unsure{trop} envie d'être de Poincaré\ldots

On peut tensoriser par $\mathbb{Q}$\ldots

3) On peut prendre $K(X)$ [\unsure{jaugé} par $\chi$], considéré comme gradué de degré 0\ldots

\figure{\begin{tikzpicture}\draw (0,0)--(5.7,0);\end{tikzpicture}}
Formalisme des correspondances. Soit $C\in A(I\amalg I')$. On va lui associer une application
\[
x\unsure{\mapstochar\relbar|\mapsto} C(x)=C.x:\quad A(I)\longrightarrow A(I')
\]
Ce sera le composé (où $\alpha:I-I\amalg I'$, $\alpha':I'-I\amalg I'$)
\diagram{
\begin{tikzcd}
A(I) \arrow[r,"\unsure{\alpha^*}"] & A(I\amalg I') \arrow[r,"y\mapsto Cy"] & A(I\amalg I') \arrow[r,"\alpha'_*"] & A(I')
\end{tikzcd}
}
\margin{il y a un choix: prendre $Cy$ ou $yC$ !}

Ces actions sont connues quand on connaît, en plus des lois d'\uline{anneaux} des $A(\unsure{\cdot})$, les $\alpha^*$ et $\alpha_*$ pour $\alpha:I\unsure{\lhook\joinrel-|\hookrightarrow}J$ une injection [il suffit avec $\operatorname{card}(J)=\operatorname{card}(I)+1$, i.e $\alpha$ un ``opérateur face''].
\margin{43}
